% !LW recipe=uplatex
\documentclass[12pt,dvipdfmx]{article}
\usepackage[margin=1.0truein]{geometry}
\usepackage{ifdraft}
\usepackage{subfiles}

% Math
\usepackage{amsmath}
\usepackage{amssymb}
\usepackage{bm}
\usepackage{mathrsfs}
\usepackage{mathtools}
\usepackage{optidef}
\usepackage{physics}
% https://qiita.com/rityo_masu/items/efd44bc8f9229e014237
\allowdisplaybreaks[4]
\DeclarePairedDelimiter{\floor}{\lfloor}{\rfloor}
\DeclarePairedDelimiter{\ceil}{\lceil}{\rceil}
% https://tex.stackexchange.com/questions/28836/typesetting-the-define-equals-symbol
% https://tex.stackexchange.com/questions/5502/how-to-get-a-mid-binary-relation-that-grows
\newcommand{\relmiddle}[1]{\mathrel{}\middle#1\mathrel{}}

% Theorem environments
\usepackage{amsthm}
\usepackage{thm-restate}
\theoremstyle{plain}
\newtheorem{theorem}{Theorem}
\newtheorem{proposition}{Proposition}
\newtheorem{lemma}{Lemma}
\newtheorem{assumption}{Assumption}
\newtheorem{condition}{Condition}
\newtheorem{corollary}{Corollary}

% Figure, Table, Algorithm
\usepackage{graphicx}
\usepackage{booktabs}
\usepackage{multirow}
\usepackage{makecell}
% \usepackage{mathptmx}
\usepackage{subcaption}
\usepackage{float}
\usepackage{tikz}
\definecolor{cA}{HTML}{0072BD}
\definecolor{cB}{HTML}{EDB120}
\definecolor{cC}{HTML}{77AC30}
\definecolor{cD}{HTML}{D95319}
\newcommand{\red}[1]{\textcolor{red}{#1}}
\newcommand{\blue}[1]{\textcolor{blue}{#1}}
\newcommand{\cyan}[1]{\textcolor{cyan}{#1}}
\newcommand{\gray}[1]{\textcolor{gray}{#1}}
\newcommand{\green}[1]{\textcolor{green}{#1}}
\newcommand{\brown}[1]{\textcolor{brown}{#1}}
\newcommand{\black}[1]{\textcolor{black}{#1}}
\newcommand{\orange}[1]{\textcolor{orange}{#1}}

% References
\usepackage{hyperref}
\usepackage{cleveref}
% cleveref fails to accurately capture pseudo-code line numbers. We define a custom reference command for lines in algorithms.
\crefname{assumption}{Assumption}{Assumptions}
\newcommand{\refLine}[1]{Line~\ref{#1}}
\newcommand{\refLines}[2]{Lines~\ref{#1} and \ref{#2}}

% Other packages
\usepackage{appendix}
\usepackage{cancel}
\usepackage{enumitem}
\usepackage{fancybox}
\usepackage{ifthen}
\usepackage{lipsum}
\usepackage{pifont}
\usepackage{xparse}

% https://www.reddit.com/r/LaTeX/comments/111e7gn/can_someone_help_me_replicate_these_boxes_in_latex/
\usepackage{tcolorbox}
\tcbuselibrary{skins, breakable}
\NewTotalTColorBox[auto counter]{\bookBox}{ +m +m }{
    notitle,
    colback=white,
    frame hidden,
    boxrule=0pt,
    enhanced,
    sharp corners,
    borderline west={4pt}{0pt}{cA},
    breakable,
}{
    \textcolor{cA}{
        \sffamily\textbf{#1}
    }#2
}
\NewTotalTColorBox[auto counter]{\myBox}{ +m +m }{
    notitle,
    colback=white,
    frame hidden,
    boxrule=0pt,
    enhanced,
    sharp corners,
    borderline west={4pt}{0pt}{cC},
    breakable,
}{
    \textcolor{cC}{
        \sffamily\textbf{#1}
    }#2
}

\begin{document}

\title{Supplemental Material for 06/30\\Book Reading Seminar}
\author{Hiroki Hamaguchi}
\date{2026/06/30}
\maketitle

\begin{abstract}
    \begin{center}
        This is the supplemental material for the book reading seminar.

        If necessary, please also refer to my \href{https://github.com/hari64boli64/BookReadingSeminarMaterials}{repository}.

        The textbook link of Springer is \href{https://link.springer.com/book/10.1007/b97543}{here}.
    \end{center}
\end{abstract}

\section*{第5章 感度解析と安定性（pp.419--445）}

\bookBox{第5章の目的と二種類の感度解析}{
    変分不等式 $\operatorname{VI}(K,F)$ の解が、集合 $K$ や写像 $F$ の小さな変化に対して
    どのように変化するかを調べる。厳密には、閉凸集合の族および連続写像の空間に距離を導入して
    「小さな変化」を定式化する必要がある。

    本章では二種類の問題を区別する。第一は、孤立解 $x^*$ を一つ固定し、摂動後の問題
    $\operatorname{VI}(L,G)$ に $x^*$ の近くの解が存在するか、その解が孤立しているか、
    摂動が消えるとき近傍の解が $x^*$ に収束するかを問う\textbf{孤立感度解析}である。
    第二は解集合全体 $\operatorname{SOL}(K,F)$ の変化を調べる\textbf{全感度解析}であり、
    こちらの方が一般に難しい。本章では発展の進んでいる孤立感度解析を主に扱う。

    パラメータ付き問題
    \[
        \{\operatorname{VI}(K(p),F(\cdot,p)):p\in P\}
    \]
    では、$p^*$ における解 $x^*$ から出る解軌道 $x(p)$ の連続性に加え、微分可能性も問題となる。
    とくに本章前半では $K$ を固定し $F$ のみを摂動する。この設定は MiCP や NCP の感度解析を含む。
}

\myBox{最適化問題との対応}{
    $F=\nabla f$ のとき、$\operatorname{VI}(K,F)$ は制約付き最適化の一階条件である。
    したがって孤立感度解析は、目的関数やパラメータを少し変えた際に局所解が残るか、
    また解がどの程度動くかを調べる問題と読める。ただし一般の VI ではポテンシャル $f$ が存在するとは限らず、
    以下の議論は勾配構造を仮定しない。
}

\bookBox{5.1 孤立解の感度:局所一意性だけでは足りない}{
    $K$ を閉凸集合として固定し、$x^*$ を $\operatorname{VI}(K,F)$ の局所一意解とする。
    中心的な問いは、$F$ に近い $G$ に対して $\operatorname{VI}(K,G)$ が $x^*$ の近くに解をもつか、である。
    しかし局所一意性だけではこれを保証できない。

    例5.1.1では
    \[
      q=(0,-1,1)^\top,\qquad
      M=\begin{bmatrix}0&-1&-1\\1&0&0\\-1&0&0\end{bmatrix}
    \]
    による LCP は一意解 $(1,0,0)$ をもつ一方、$q(\varepsilon)=(-\varepsilon,-1,1)^\top$
    とすると任意の $\varepsilon>0$ で解が存在しない。$M$ は半正定値ではあるが正定値でない。
    ゆえに「元の解が孤立していること」と「近い問題に近い解が存在すること」は別の性質である。

    開集合 $D\subseteq\mathbb R^n$ 上の連続写像 $H$ と $S\subset D$ に対し、
    \[
       \mathbb B(H;\varepsilon,S)
       \coloneqq\left\{G:\sup_{y\in S}\norm{G(y)-H(y)}<\varepsilon\right\}
    \]
    を $S$ 上の一様な $\varepsilon$-近傍とする。$\norm{G_k-H}_S\to0$ を
    $G_k$ が $S$ 上で $H$ に収束することと呼ぶ。
}

\bookBox{命題5.1.2:孤立解は近傍解の吸引点である}{
    $F\colon D\to\mathbb R^n$ を連続とし、$x^*$ を $\operatorname{VI}(K,F)$ の孤立解とする。
    $x^*$ の近傍 $N$ が
    \[
       \overline N\subset D,\qquad
       \operatorname{SOL}(K,F)\cap\overline N=\{x^*\}                         \tag{5.1.1}
    \]
    を満たすとする。$G_k\to F$ が $N$ 上で一様に成り立ち、
    $x_k\in\operatorname{SOL}(K,G_k)\cap N$ ならば $x_k\to x^*$ である。

    証明の要点はコンパクト性と極限操作である。$\{x_k\}$ の任意の集積点 $x^\infty$ は、
    $G_k\to F$ と VI の不等式を極限に移すことにより $\operatorname{VI}(K,F)$ の解となる。
    しかも $x^\infty\in\overline N$ なので (5.1.1) より $x^\infty=x^*$ である。
    集積点が一つしかないため列全体が $x^*$ に収束する。

    ただしこの命題は、摂動後の解が近傍に\emph{存在すると仮定した上で}その収束を述べるだけであり、
    存在そのものは保証しない。局所可解性には位相的次数を用いる。
}

\bookBox{命題5.1.3--5.1.5:非零指数による局所可解性}{
    $F_K^{\mathrm{nat}}$ を自然写像とする。孤立解 $x^*$ における指数
    $\operatorname{ind}(F_K^{\mathrm{nat}},x^*)$ が非零ならば、(5.1.1) を満たす任意の開近傍 $N$ に対し、
    \[
       \varepsilon=\operatorname{dist}_\infty
       (0,F_K^{\mathrm{nat}}(\partial N))>0
    \]
    とおくことで、$\norm{G-F}_{\overline N}<\varepsilon$ を満たす $G$ による
    $\operatorname{VI}(K,G)$ は $N$ 内に解をもつ（命題5.1.3）。射影 $\Pi_K$ の非拡大性により
    自然写像同士の差もこの誤差以下となり、次数の近接不変性から
    \[
       \deg(G_K^{\mathrm{nat}},N)
       =\operatorname{ind}(F_K^{\mathrm{nat}},x^*)\ne0
    \]
    が従うためである。

    正規写像を用いると仮定を局所化できる。$z^*=x^*-F(x^*)$、$K_N=K\cap\overline N$ とおく。
    $\operatorname{ind}(F_K^{\mathrm{nor}},z^*)\ne0$ ならば、ある $\varepsilon>0$ が存在し、
    $G-F$ の大きさを $K_N$ 上だけで抑えれば $\operatorname{VI}(K,G)$ は $N$ に解をもつ
    （命題5.1.4）。これは $\Pi_K(z)$ が $K_N$ に入る正規写像側の近傍を選べるためである。

    パラメータ付きの固定集合 VI では、$F(\cdot,p^*)$ に対する正規写像の指数が非零で、
    \[
       \sup_{x\in K\cap N}\norm{F(x,p)-F(x,p^*)}
       \le L\norm{p-p^*}
    \]
    が成り立つなら、$p^*$ の十分小さい近傍で
    $S_N(p)=\operatorname{SOL}(K,F(\cdot,p))\cap N$ は空でない。さらに
    \[
       \lim_{p\to p^*}\sup\{\norm{x-x^*}:x\in S_N(p)\}=0.                 \tag{5.1.3}
    \]
    これは命題5.1.4で存在を、命題5.1.2で収束を得る二段構えである。
}

\myBox{位相的次数の役割}{
    非零次数は、境界を通って零点が逃げない限り内部の零点が消滅できないことを表す。
    したがってここでの指数条件は「解の個数」を直接数える条件というより、
    摂動下で少なくとも一つの解を近傍に残すための位相的な証明書と理解するとよい。
    一意性と存在性を別々に扱う本節の構成も、この役割分担を反映している。
}

\bookBox{命題5.1.6--系5.1.8:臨界錐上の正値性と定量評価}{
    $F$ が解 $x^*$ で B微分可能とする。臨界錐
    \[
       C(x^*;K,F)=T(x^*;K)\cap F(x^*)^\perp
    \]
    上で
    \[
       v^\top F'(x^*;v)>0
       \qquad(0\ne v\in C(x^*;K,F))                                  \tag{5.1.4}
    \]
    が成り立つなら、自然写像と正規写像の指数はいずれも定義され非零となる（命題5.1.6）。
    証明では $F$ と $x\mapsto x-z^*$ を結ぶホモトピー
    \[
      F_t(x)=tF(x)+(1-t)(x-z^*)
    \]
    を作る。もしホモトピーの零点が任意に小さい近傍の境界へ現れるなら、正規化した解差の極限から
    非零臨界方向 $v$ が得られ、(5.1.4) に反する。ゆえに次数は恒等写像の場合と同じ $1$ である。

    同じ条件は解の移動量も制御する。任意の適切な近傍 $N$ に対してある $c>0$ が存在し、
    \[
      \sup\{\norm{x-x^*}:x\in\operatorname{SOL}(K,q+F)\cap\overline N\}
      \le c\norm{q}       \tag{5.1.7}
    \]
    となる（補題5.1.7）。これと局所可解性を合わせると、ある $c,\varepsilon>0$ が存在して、
    \[
      \gamma\coloneqq\sup_{x\in K\cap\overline N}\norm{G(x)-F(x)}<\varepsilon
    \]
    ならば $\operatorname{SOL}(K,G)\cap N\ne\varnothing$ かつ
    \[
      \sup_{x\in\operatorname{SOL}(K,G)\cap\overline N}\norm{x-x^*}
      \le c\gamma.                                                        \tag{5.1.5}
    \]
    パラメータ付きの場合には $\gamma\le L\norm{p-p^*}$ より、局所解写像について
    上側 Lipschitz 性が得られる。
}

\myBox{二次十分条件との類似}{
    (5.1.4) は、最適化において臨界方向上で曲率が正であるという二次十分条件に対応する形をしている。
    $F=\nabla f$ が滑らかなら $F'(x^*;v)=\nabla^2f(x^*)v$ なので、左辺は
    $v^\top\nabla^2f(x^*)v$ となる。この正値性が局所一意性だけでなく、摂動解の存在と
    $O(\text{摂動量})$ の誤差評価まで与える点が重要である。
}

\bookBox{5.2 B微分可能方程式:半安定性と安定性}{
    $K=\mathbb R^n$ とすれば VI は方程式 $H(x)=0$ になる。$H$ が零点 $x$ で B微分可能なとき、
    次の三条件は同値である（補題5.2.1）:
    \begin{enumerate}[label=(\alph*)]
      \item $H'(x;v)=0$ ならば $v=0$;
      \item ある $c'>0$ により $\norm{v} \le c'\norm{H'(x;v)}$ が全ての $v$ で成り立つ;
      \item ある近傍 $N$ と $c>0$ により $\norm{y-x}\le c\norm{H(y)}$ が全ての $y\in N$ で成り立つ。
    \end{enumerate}
    とくにこれらは $x$ の孤立性を含意する。さらに
    $\Gamma(y)=H'(x;y-x)$ とおいて $\operatorname{ind}(\Gamma,x)\ne0$ なら、
    十分小さい一様摂動 $G$ は $x$ の近傍に零点をもつ（補題5.2.2）。B微分が $H$ の一次近似であるため
    $\operatorname{ind}(\Gamma,x)=\operatorname{ind}(H,x)$ となり、次数論を適用できる。

    これを踏まえ、孤立零点 $x$ を次のように定義する（定義5.2.3）。
    \begin{itemize}
      \item \textbf{半安定}:小さい任意の摂動 $G$ の近傍零点 $x'$ が
      $\norm{x'-x}\le c\norm{H(x')}$ を満たす。ただし摂動後の零点の存在は要求しない。
      \item \textbf{安定}:上の誤差評価に加えて、全ての十分小さい摂動 $G$ が近傍に零点をもつ。
    \end{itemize}
    よって $H'(x;\cdot)$ の零点が原点だけなら半安定であり、さらに指数が非零なら安定である
    （定理5.2.4）。これらは局所概念であり、十分小さな近傍だけを考えればよい。
}

\bookBox{強安定性と正則性（定義5.2.6--5.2.7）}{
    任意の連続摂動に対して零点の一意性まで要求するのは強すぎる。実際 $H(t)=t$ でさえ、
    $G(t)=t-\varepsilon'\sqrt{|t|}\sin(1/\sqrt{|t|})$ のような一様に小さい連続摂動を選べば、
    任意の零点近傍に複数の零点を作れる（例5.2.5）。

    そこで孤立零点 $x$ が\textbf{強安定}であるとは、十分小さい任意の摂動 $G,\widetilde G$ が近傍零点をもち、
    それらの任意の零点 $v,\widetilde v$ に対して
    \[
       \norm{v-\widetilde v}
       \le c\norm{H(v)-H(\widetilde v)}
    \]
    が成り立つことと定める。$\widetilde G=H$ とすれば半安定性が従うので、強安定なら安定である。
    とくに右辺摂動 $H(y)=u$ に限れば、近傍解 $x(u)$ は一意で
    \[
       \norm{x(u)-x(u')}\le c\norm{u-u'}
    \]
    を満たす。

    右辺摂動だけに限定した概念が\textbf{正則性}である。正則なら
    $\varnothing\ne H^{-1}(u)\cap N\subset \mathbb B(x,c\norm{u})$、
    \textbf{強正則}ならさらにこの集合は一元集合 $\{x_N(u)\}$ で、$x_N$ は Lipschitz 連続である。
    任意の関数摂動を許す安定性は、対応する正則性より強い。
}

\bookBox{定理5.2.8--命題5.2.9:強安定性の特徴付け}{
    連続写像 $H$ の零点 $x$ について、強安定性と強正則性は同値である（定理5.2.8(a)）。
    強正則性から強安定性を示す際には、局所逆写像 $x_N$ と誤差 $e=H-G$ を使って
    \[
       \Phi(y)=x_N(e(y))
    \]
    を閉球からそれ自身への連続写像として構成し、Brouwer の不動点定理を適用する。

    また $H$ が $x$ で局所 Lipschitz 同相なら $x$ は強安定である。逆に $H$ 自身が $x$ で
    局所 Lipschitz 連続なら、強安定性から局所 Lipschitz 同相性が従う。したがって局所 Lipschitz 写像では
    \[
       \text{強安定}\quad\Longleftrightarrow\quad
       \text{強正則}\quad\Longleftrightarrow\quad
       \text{局所 Lipschitz 同相}.
    \]

    さらに $H$ が $x$ で強 F微分可能なら、次の全てが同値である（命題5.2.9）:
    \[
      \text{強安定},\quad \text{安定},\quad \text{正則},\quad
      JH(x)\text{ が正則},\quad \text{局所 Lipschitz 同相}.
    \]
    滑らかな場合には結局、古典的な逆関数定理の非特異 Jacobian 条件へ帰着する。
}

\bookBox{垂直相補性問題と非滑らかな反例}{
    $H(x)=\min(F(x),G(x))$ とおけば、$H(x)=0$ は垂直相補性問題
    \[
       0\le F(x)\perp G(x)\ge0
    \]
    と同値である。解 $x^*$ において
    \[
      \alpha_F=\{i:F_i(x^*)=0<G_i(x^*)\},\quad
      \beta=\{i:F_i(x^*)=G_i(x^*)=0\},\quad
      \alpha_G=\{i:F_i(x^*)>0=G_i(x^*)\}
    \]
    とする。非退化性 $\beta=\varnothing$ の下では、$H$ の各行は近傍で $F_i$ または $G_i$ の一方に固定される。
    したがって
    \[
       \begin{bmatrix}
          \nabla F_i(x^*)^\top:i\in\alpha_F\\
          \nabla G_i(x^*)^\top:i\in\alpha_G
       \end{bmatrix}
    \]
    が正則なら、十分小さい二組の摂動問題は近傍に一意解 $x_1,x_2$ をもち、解差は摂動差で
    Lipschitz 評価できる（系5.2.10）。

    一方、非滑らかな点では安定性と強安定性は一致しない。例5.2.11の
    \[
       H(x_1,x_2)=
       \begin{bmatrix}\min(x_1,x_1+x_2)\\\min(x_2,x_1+x_2)\end{bmatrix}
    \]
    は原点を一意零点にもち、対応 LCP の厳密共正値性と指数 $1$ から安定である。
    しかし $H(x)+(\varepsilon,\varepsilon)^\top=0$ は
    $(-\varepsilon,0)$ と $(0,-\varepsilon)$ の少なくとも二解をもつため強安定ではない。
    この例は、退化 VI・CP を扱うには滑らかな逆関数定理だけでは不十分であることを示す。
}

\bookBox{定理5.2.12・5.2.14:B微分による必要十分条件}{
    $\Gamma(y)=H'(x;y-x)$ とする。$\Gamma$ が $x$ で\textbf{非消失性}をもつとは、ある連続写像 $\Phi$ が存在し、
    全ての $\delta\ge0$ について $x$ が $\Phi+\delta\Gamma$ の唯一の零点であり、かつ
    $\operatorname{ind}(\Phi,x)\ne0$ となることをいう。正則な線形部をもつアフィン写像や
    $P_0$ 写像はこの性質をもつ。

    $H$ が $x$ で B微分可能で $\Gamma$ が非消失性をもつなら、次は同値である（定理5.2.12）:
    \begin{enumerate}[label=(\alph*)]
      \item $x$ は $\Gamma$ の安定零点である;
      \item $\Gamma$ は全射で、$\Gamma^{-1}(\mathbb B(0,1))$ は有界である;
      \item $x$ は $\Gamma$ の唯一の零点である;
      \item $x$ は $H$ の安定零点である。
    \end{enumerate}
    さらに安定性の定義で「全ての近傍」を要求する代わりに「ある一つの近傍」で条件が成り立てば十分である。

    $H$ が $x$ で\textbf{強 B微分可能}なら、非消失性の仮定なしに次が同値となる（定理5.2.14）:
    \[
       0\text{ は }H'(x;\cdot)\text{ の強安定零点}
       \Longleftrightarrow H'(x;\cdot)\text{ は大域的 Lipschitz 同相}
       \Longleftrightarrow x\text{ は }H\text{ の強安定零点}.
    \]
    B微分が区分線形なら、第4章の区分線形写像に対する大域的同相条件も利用できる。
}

\bookBox{5.2.2 一次近似（FOA）への拡張}{
    局所 Lipschitz 写像 $H$ に対し、$h$ が $x$ における一次近似（FOA）であるとは
    \[
       \lim_{y\to x,\,y\ne x}\frac{\norm{H(y)-h(y)}}{\norm{y-x}}=0
    \]
    を満たすことをいう。誤差 $e=H-h$ が
    \[
       \lim_{(y_1,y_2)\to(x,x),\,y_1\ne y_2}
       \frac{\norm{e(y_1)-e(y_2)}}{\norm{y_1-y_2}}=0
    \]
    を満たすとき強 FOA と呼ぶ。FOA であることは $H$ と $h$ について対称である。
    $H$ が B微分可能なら $h(y)=H(x)+H'(x;y-x)$ が FOA となる。

    正規写像 $F_K^{\mathrm{nor}}$ では、射影 $\Pi_K$ を線形化せず $F$ だけを線形化した
    半線形化 VI の正規写像が FOA になる。すなわち $\bar x=\Pi_K(\bar z)$ として
    \[
      z\mapsto F_K^{\mathrm{nor}}(\bar z)
      +JF(\bar x)(\Pi_K(z)-\Pi_K(\bar z))+z-\Pi_K(z).                 \tag{5.2.8}
    \]
    これは射影まで近似する B微分近似とは異なる。

    命題5.2.15によれば、$h$ が $H$ の FOA で、$x$ が $H$ の安定零点かつ
    $\operatorname{ind}(H,x)\ne0$ なら、$x$ は $h$ の安定零点である。
    また $h$ が強 FOA なら、$x$ が $H$ の強安定零点であることと $h$ の強安定零点であることは同値である。
    前半で指数条件を除けるかは明らかでない（注意5.2.16）。
}

\myBox{この範囲の見取り図}{
    本範囲の条件は、おおむね次の三段階を担う。
    \[
      \begin{array}{c}
      \text{方向微分の零点が原点のみ}\\[-1mm]
      \Downarrow\\[-1mm]
      \text{誤差評価・孤立性（半安定）}
      \end{array}
      \qquad
      \begin{array}{c}
      +\ \text{非零指数}\\[-1mm]
      \Downarrow\\[-1mm]
      \text{摂動解の存在（安定）}
      \end{array}
      \qquad
      \begin{array}{c}
      +\ \text{局所逆写像の一意性}\\[-1mm]
      \Downarrow\\[-1mm]
      \text{Lipschitz な解応答（強安定）}.
      \end{array}
    \]
    滑らかな問題ではこれらが Jacobian の正則性へまとまるが、非滑らか・退化した VI では分離する。
    次節5.3は、この方程式に対する概念を自然写像・正規写像を介して固定集合 VI に移し、
    VI 固有の構造を反映した安定性を定義するところから始まる。
}

\end{document}
